%==============================================================
%------------------JONATHAN TOBIAS DA SILVA--------------------
%--------------------ENGENHARIA ELÉTRICA-----------------------
%------------INSTITUTO FEDERAL- CAMPUS SERTÃOZINHO-------------
%--------------main.tex, v1.0.0, JonathanTSilva----------------
%
%------------------github.com/JonathanTSilva-------------------
%---------------linkedin.com/in/JonathanTSilva-----------------
%
% Copyright, 2012–2021 by Abntex2 at https://www.abntex.net.br/
%==============================================================

\documentclass[
	%---------------OPÇÕES DA CLASSE MEMOIR-------------------|
	12pt,				% tamanho da fonte
	openright,			% capítulos começam em pág ímpar (insere página vazia caso preciso)
	oneside,			% para impressão em verso e anverso. Oposto a "oneside" e "twoside"
	a4paper,			% tamanho do papel. 
	%---------------OPÇÕES DA CLASSE ABNTEX2------------------|
	chapter=TITLE,		% títulos de capítulos convertidos em letras maiúsculas
	%section=TITLE,		% títulos de seções convertidos em letras maiúsculas
	%subsection=TITLE,	% títulos de subseções convertidos em letras maiúsculas
	%subsubsection=TITLE,% títulos de subsubseções convertidos em letras maiúsculas
	%----------------OPÇÕES DO PACOTE BABEL-------------------|
	english,			% idioma adicional para hifenização
	%french,			% idioma adicional para hifenização
	%spanish,			% idioma adicional para hifenização
	brazil				% o último idioma é o principal do documento
	]{abntex2}

%-------------------------PACOTES-----------------------------|
%==============================================================
%------------------JONATHAN TOBIAS DA SILVA--------------------
%--------------------ENGENHARIA ELÉTRICA-----------------------
%------------INSTITUTO FEDERAL- CAMPUS SERTÃOZINHO-------------
%-------------packages.tex, v1.0.0, jonathanTSilva-------------
%
%------------------github.com/JonathanTSilva-------------------
%---------------linkedin.com/in/JonathanTSilva-----------------
%
% Copyright, 2012–2021 by abnTeX2 at https://www.abntex.net.br/
%==============================================================

%--------------------LÍNGUA E CARACTERES:---------------------|
\usepackage[utf8]{inputenc}                                     % Codificação do documento (conversão automática dos acentos)
\usepackage[T1]{fontenc}                                        % Seleção de códigos de fonte
%\usepackage{fontspec}                                          % Seleção de fontes
%-------------------------------------------------------------|

%-----------------------MODO MATEMÁTICO-----------------------|
\usepackage{amsmath, amsfonts, amssymb}                         % Pacotes da sociedade americana de matemática ASM
%-------------------------------------------------------------|

%--------------------------CRIAÇÕES---------------------------|
%\usepackage{tikz}
%-------------------------------------------------------------|

%---------------------EXIBIÇÃO E FORMATAÇÃO-------------------|
\usepackage{graphicx}                                           % Inclusão de gráficos
\usepackage[table, svgname]{xcolor}                             % Controle de Cores
\usepackage{colortbl}
\usepackage{indentfirst}                                        % Indenta o primeiro parágrafo de cada seção
\usepackage{float}                                              % Fixar imagem no ponto específicado [H]
\usepackage{multirow, multicol}                                 % Tabelas com células multicolunas e multilinhas
\usepackage{enumitem}                                           % Controle sobre o layout dos três ambientes básicos de lista: enumerar, itemizar e descrição
\usepackage{microtype}                                          % Para melhorias de justificação
\usepackage{listings}                                           % Inserir código de linguagem de programação
\usepackage{nomencl}                                            % Necessário para o commando makeindex
%\usepackage{floatrow}
\usepackage{subfig}
\usepackage[size=footnotesize, skip=3pt]{caption}
%\usepackage{subcaption}
\usepackage{longtable}                                          % Para as tabelas
\usepackage{booktabs}
\usepackage{tabularx}
%-------------------------------------------------------------|

%---------------BIBLIOGRAFIA, LINKS E CITAÇÕES----------------|
\usepackage[brazilian, hyperpageref]{backref}	                % Paginas com as citações na bibl
\usepackage{xurl}
\usepackage{hyperref}                                           % Três pacotes necessários para a quebra de url longa
\usepackage[hyphenbreaks]{breakurl}
\usepackage[alf,abnt-emphasize=bf]{abntex2cite}	                % Citações padrão ABNT, quando classe abnTeX2
%\citeoption{minhasopcoes}                                      % Le as opções estabelecidas no arquivo .bib para abnTeX2cite
%-------------------------------------------------------------|

%---------------------------OUTROS----------------------------|
\usepackage{pdfpages}                                           % Inclui página de outro pdf. 
                                                                % Ex: \includepdf[pages={1,3-},nup=1x2,landscape=true]{main.pdf}
\usepackage{lastpage}                                           % Usado pela Ficha catalográfica
%-------------------------------------------------------------|

%---------------------------FONTES----------------------------|
\usepackage{lmodern}                                            % Usa a fonte Latin Modern
%-------------------------------------------------------------|

%--------------------------FONTES-----------------------------|
%-----ARIAL---------------%
%\usepackage{helvet}
%\renewcommand{\familydefault}{\sfdefault}

%-----TIMES NEW ROMAN-----%
%\usepackage{times}

%-----ADOBE UTOPIA--------%
%\usepackage{fourier}

%-------------------CAPA E FOLHA DE ROSTO---------------------|
%==============================================================
%--------------------CAPA E FOLHA DE ROSTO---------------------
%--------------main.tex, v1.0.0, JonathanTSilva----------------
%==============================================================

\titulo{[TÍTULO DO TCC]}
\autor{[NOME DO AUTOR 1] \and [NOME DO AUTOR 2]}
\local{Sertãozinho}
\data{2021}
\orientador{[ORIENTADOR]}
\coorientador{[COORIENTADOR]}
\instituicao{%
    Instituto Federal de Educação, Ciência e Tecnologia de São Paulo - IFSP \par
    [BACHARELADO OU LICENCIATURA] em [CURSO] \par
    Trabalho de Conclusão de Curso
}
\tipotrabalho{Monografia (Trabalho de Conclusão de Curso)}
\preambulo{Trabalho de Conclusão de Curso apresentado à
Coordenação do Curso de [NOME DO CURSO] do Instituto Federal de Educação, Ciência e Tecnologia de São Paulo - IFSP - Campus Sertãozinho, como requisito parcial para obtenção do título de [BACHARELADO OU LICENCIATURA] em [CURSO].} 

%==============================================================
%----------------------CONFIGURAÇÕES---------------------------
%==============================================================

%----------------------PACOTE: HYPEREF------------------------|
\makeatletter
\hypersetup{
    pdftitle = {\@title},
    pdfauthor = {\@author},
    pdfsubject = {\imprimirpreambulo},
    pdfborder = {0 0 0},
    pdfkeywords = {}{}{}{},
    pdfcreator = {LaTeX with abnTeX2},
    colorlinks = true,
    linkcolor = blue,
    citecolor = blue,
    urlcolor = blue
}
\makeatother

%----------------------PACOTE: BACKREF------------------------|
\renewcommand{\backrefpagesname}{Citado na(s) página(s):~}  % Usado sem a opção hyperpageref de backref
\renewcommand{\backref}{}                                   % Texto padrão antes do número das páginas
\renewcommand*{\backrefalt}[4]{                             % Define os textos da citação
	\ifcase #1%
		Nenhuma citação no texto.%
	\or
		Citado na página #2.%
	\else
		Citado #1 vezes nas páginas #2.%
	\fi}%

%-------------------------APARÊNCIA---------------------------|
\definecolor{blue}{RGB}{41,5,195}
\definecolor{gray}{RGB}{190,190,190}

%--------------------------TABELAS----------------------------|
\newcommand*{\thead}[1]{\multicolumn{1}{c}{\bfseries #1}}

%------POSICIONAMENTO DE FIGS E TABS EM PÁGINA EM BRANCO------|
% Posiciona figuras e tabelas no topo da página quando adicionadas sozinhas
% em um página em branco. Ver https://github.com/abntex/abntex2/issues/170
\makeatletter
\setlength{\@fptop}{5pt} % Define a distância de topo da página no primeiro flutuante
\makeatother

\iffalse %COMENTADO
    %---------CRIAÇÃO DE QUADROS E LISTA DE QUADROS-----------|
    % Possibilita criação de Quadros e Lista de quadros.
    % Ver https://github.com/abntex/abntex2/issues/176
    \newcommand{\quadroname}{Quadro}
    \newcommand{\listofquadrosname}{Lista de quadros}
    
    \newfloat[chapter]{quadro}{loq}{\quadroname}
    \newlistof{listofquadros}{loq}{\listofquadrosname}
    \newlistentry{quadro}{loq}{0}
    
    % configurações para atender às regras da ABNT
    \setfloatadjustment{quadro}{\centering}
    \counterwithout{quadro}{chapter}
    \renewcommand{\cftquadroname}{\quadroname\space} 
    \renewcommand*{\cftquadroaftersnum}{\hfill--\hfill}
    
    \setfloatlocations{quadro}{hbtp}
    
    %\floatsetup[figure]{capposition=top}
\fi

%------------------------ESPAÇAMENTO--------------------------|
\setlength{\parindent}{1.3cm} % O tamanho do parágrafo
\setlength{\parskip}{0.2cm}   % tente também \onelineskip; Espaçamento entre um par e outro
\OnehalfSpacing

%--------------------------TÍTULOS----------------------------|
%\renewcommand{\ABNTEXchapterfont}{\fontfamily{cmr}\fontseries{b}\selectfont}
%\renewcommand{\ABNTEXchapterfont}{\bfseries} % Deixar título de capítulos em negrito

%----------------------COMPILA O ÍNDICE-----------------------|
\makeindex


%==============================================================
%-------------------------DOCUMENTO----------------------------
%==============================================================

\begin{document}
    %------------------------NOTA-----------------------------|
    %=======================EXCLUIR============================
    \label{sec:README}
    %==============================================================
%------------------JONATHAN TOBIAS DA SILVA--------------------
%---------------ENGENHARIA ELÉTRICA - 1ª TURMA-----------------
%-----------INSTITUTO FEDERAL - CAMPUS SERTÃOZINHO-------------
%--------------Main.tex, v1.0.0, JonathanTSilva----------------
%==============================================================

\documentclass[../Main.tex]{subfiles}

\begin{document}
    \begin{minipage}{0.23\linewidth}
    	\begin{figure}[H]
    		\includegraphics[scale=0.30]{Images/logo-ifsp}
    	\end{figure}
    \end{minipage}
    \begin{minipage}{0.39\linewidth}
    	\begin{table}[H]
    	\def\arraystretch{1} %esp. entre linhas
    		\begin{tabular}{c}
    			\textbf{JONATHAN TOBIAS DA SILVA} \\
    			jonathantobias2009@hotmail.com \\
    			Engenharia Elétrica \\
    			Instituto Federal de São Paulo
    		\end{tabular}
    	\end{table}
    \end{minipage}
    \begin{minipage}{0.26\linewidth}
    	\begin{figure}[H]
    		\flushright
    		\includegraphics[scale=0.07]{Images/logo-eng2}
    	\end{figure}
    \end{minipage}
    
    \vspace{0.5cm}
    
    \mysection{NOTAS}
        \begin{enumerate}
            \item {\color{red} ESTA SEÇÃO DEVERÁ SER EXCLUÍDA DO SEU DOCUMENTO FINAL.}
            
            \item Este documento foi desenvolvido em \LaTeX\ como \textit{template} do Relatório Parcial necessário para o acompanhamento de um projeto de Iniciação Científica do Instituto Federal de Educação, Ciência e Tecnologia do Estado de São Paulo - IFSP - Campus Sertãozinho.
            
            \item Este documento está disponível no repositório do GitHub:
            
            \begin{center}
                \href{https://github.com/JonathanTSilva}{github.com/JonathanTSilva}
            \end{center}
            
            \item Todos os campos que deverão ser alterados estão apresentados na forma de colchetes: [CAMPO].
            
            \item Para a exclusão desta seção, ir para o arquivo principal  (Main.tex) e deletar as seguintes linhas:
            
\begin{verbatim}
%=======================EXCLUIR============================
\label{sec:README}
\subfile{Sections/README}
\newpage
%==========================================================
\end{verbatim}
        \end{enumerate}
         
    \mysection{HISTÓRICO DE MODIFICAÇÕES}
        A \autoref{tab:modificacoes} apresenta o histórico de modificações deste documento, que deverão ser anotadas e descritas as versões desenvolvidas deste \textit{template}, acrescentando na tabela o respectivo número de versão, a elaboração, aprovação e breve descrição da alteração realizada.
        
        \begin{table}[H]
            \centering
            \caption{Histórico de modificações}
            \begin{tabular}{|>{\centering\arraybackslash}m{0.07\textwidth}|>{\centering\arraybackslash}m{0.25\textwidth}|>{\centering\arraybackslash}m{0.25\textwidth}|>{\centering\arraybackslash}m{0.3\textwidth}|}
                \hline
                \textbf{V.}         & \textbf{ELABORAÇÃO}                   & \textbf{APROVAÇÃO}                    & \textbf{DESCRIÇÃO} \\ \hline
                \multirow{3}{*}{1.0.0}  & \hrulefill                            & \hrulefill                            & {\scriptsize Modelo de relatório -} \\
                                        & {\scriptsize Jonathan T. da Silva}    & {\scriptsize Jonathan T. da Silva}    & {\scriptsize Criação da \textit{template} em} \LaTeX \\
                                        & {\scriptsize 30/Jul, 2020}            & {\scriptsize 30/Jul, 2020}            & {\scriptsize } \\ \hline
                                        
                                        
                                        
                                        
                \multirow{3}{*}{}       & \hrulefill                            & \hrulefill                            & {\scriptsize [DESCRIÇÂO LINHA1]} \\
                                        & {\scriptsize [NOME DO ELABORADOR]}    & {\scriptsize [NOME DO APROVADOR]}     & {\scriptsize [DESCRIÇÂO LINHA2]} \\
                                        & {\scriptsize [DATA DA ELABORAÇÃO]}    & {\scriptsize [DATA DA APROVAÇÃO]}     & {\scriptsize [DESCRIÇÂO LINHA3]} \\ \hline
            \end{tabular}
            \label{tab:modificacoes}
        \end{table}
        
\end{document}
    \newpage
    %==========================================================
    
    %-----------------------IDIOMA----------------------------|
    %\selectlanguage{english}
    \selectlanguage{brazil}

    \frenchspacing % Retira espaço extra obsoleto entre as frases.
    
    %__________________________________________________________
    %----------------ELEMENTOS PRÉ-TEXTUAIS--------------------
    
    %------------------------CAPA-----------------------------|
    \begin{figure}[H]
        \centering
        \includegraphics[scale=0.45]{Figuras/IFSP.png}
        \label{fig:logoIFSP}
    \end{figure}
    
    \vspace{-1cm}
    
    \begin{center}
        \textbf{\sffamily\large Instituto Federal de Educação, Ciência e Tecnologia de São Paulo \\ Campus Sertãozinho}
    \end{center}
    
    \vspace{3cm}
    
    \imprimircapa
    
    %-------------------FOLHA DE ROSTO------------------------|
    \imprimirfolhaderosto*{} % (o * indica que haverá a ficha bibliográfica)
    %-----------------FICHA CATALOGRÁFICA---------------------|
    %==============================================================
%--------------------FICHA CATALOGRÁFICA-----------------------
%--------------main.tex, v1.0.0, JonathanTSilva----------------
%==============================================================

% Isto é um exemplo de Ficha Catalográfica, ou ``Dados internacionais de
% catalogação na publicação''. Pode-se utilizar esse modelo como referência. 
% Se utilizar um PDF fornecido pela universidade, optar pelo comando abaixo:

%\begin{fichacatalografica}
%    \includepdf{fig_ficha_catalografica.pdf}
%\end{fichacatalografica}

\begin{fichacatalografica}

	\vspace*{\fill}					% Posição vertical
	\hrule							% Linha horizontal
	\begin{center}					% Minipage Centralizado
    	\begin{minipage}[c]{12.5cm}	% Largura
    	
        	\imprimirautor
        	
        	\hspace{0.5cm} \imprimirtitulo  / \imprimirautor. --
        	\imprimirlocal, \imprimirdata-
        	
        	\hspace{0.5cm} \pageref{LastPage} p. : il. (algumas color.) ; 30 cm. \\
        	
        	\hspace{0.5cm} \imprimirorientadorRotulo~\imprimirorientador \\ \imprimirorientador \\
        	
        	\hspace{0.5cm}
        	\parbox[t]{\textwidth}{\imprimirtipotrabalho~--~\imprimirinstituicao,
        	\imprimirdata.}\\
        	
        	\hspace{0.5cm}
    		1. Palavra-chave1.
    		2. Palavra-chave2.
    		I. Orientador.
    		II. Universidade xxx.
    		III. Faculdade de xxx.
    		IV. Título \\ 			
        	
        	\hspace{8.75cm} CDU 02:141:005.7 \\
    	
    	\end{minipage}
	\end{center}
	\hrule

\end{fichacatalografica}

    %-----------------------ERRATA----------------------------|
    %==============================================================
%---------------------------ERRATA-----------------------------
%--------------main.tex, v1.0.0, JonathanTSilva----------------
%==============================================================

\begin{errata}

    Elemento opcional NBR14724:2011. Exemplo:
    
    \vspace{\onelineskip}
    
    FERRIGNO, C. R. A. \textbf{Tratamento de neoplasias ósseas apendiculares com
    reimplantação de enxerto ósseo autólogo autoclavado associado ao plasma
    rico em plaquetas}: estudo crítico na cirurgia de preservação de membro em
    cães. 2011. 128 f. Tese (Livre-Docência) - Faculdade de Medicina Veterinária e
    Zootecnia, Universidade de São Paulo, São Paulo, 2011.
    
    \begin{table}[H]
        \center
        \footnotesize
        \begin{tabular}{|p{1.4cm}|p{1cm}|p{3cm}|p{3cm}|}
            \hline
            \textbf{Folha}  & \textbf{Linha}    & \textbf{Onde se lê}   & \textbf{Leia-se} \\
            \hline
            1               & 10                & auto-conclavo         & autoconclavo \\
            \hline
        \end{tabular}
    \end{table}
    
\end{errata}
    %---------------------DEDICATÓRIA-------------------------|
    %==============================================================
%------------------------DEDICATÓRIA---------------------------
%--------------main.tex, v1.0.0, JonathanTSilva----------------
%==============================================================

\begin{dedicatoria}

    \vspace*{\fill}
    \centering
    \noindent
    \textit{Este trabalho é dedicado às crianças adultas que, \\
    quando pequenas, sonharam em se tornar cientistas.} \vspace*{\fill}
    
\end{dedicatoria}
    %-------------------AGRADECIMENTOS------------------------|
    %==============================================================
%-----------------------AGRADECIMENTOS-------------------------
%--------------main.tex, v1.0.0, JonathanTSilva----------------
%==============================================================

\begin{agradecimentos}

    Os agradecimentos principais são direcionados à ...

    \iffalse
        Leslie H. Watter, Bruno Parente Lima, Flávio de Vasconcellos Corrêa, Otavio Real
        Salvador, Renato Machnievscz\footnote{Os nomes dos integrantes do primeiro
        projeto abn\TeX\ foram extraídos de
        \url{http://codigolivre.org.br/projects/abntex/}} e todos aqueles que
        contribuíram para que a produção de trabalhos acadêmicos conforme
        as normas ABNT com \LaTeX\ fosse possível.
        
        Agradecimentos especiais são direcionados ao Centro de Pesquisa em Arquitetura
        da Informação\footnote{\url{http://www.cpai.unb.br/}} da Universidade de
        Brasília (CPAI), ao grupo de usuários
        \emph{latex-br}\footnote{\url{http://groups.google.com/group/latex-br}} e aos
        novos voluntários do grupo
        \emph{\abnTeX}\footnote{\url{http://groups.google.com/group/abntex2} e
        \url{http://abntex2.googlecode.com/}}~que contribuíram e que ainda
        contribuirão para a evolução do \abnTeX.
    \fi
    
\end{agradecimentos}
    %----------------------EPÍGRAFE---------------------------|
    %==============================================================
%--------------------------EPÍGRAFE----------------------------
%--------------main.tex, v1.0.0, JonathanTSilva----------------
%==============================================================

\begin{epigrafe}

    \vspace*{\fill}
	
	\begin{flushright}
	    \textit{``Não vos amoldeis às estruturas deste mundo, \\
		mas transformai-vos pela renovação da mente, \\
		a fim de distinguir qual é a vontade de Deus: \\
		o que é bom, o que Lhe é agradável, o que é perfeito.\\
		(Bíblia Sagrada, Romanos 12, 2)}
	\end{flushright}
	
\end{epigrafe}

    %----------------------RESUMOS----------------------------|
    %==============================================================
%---------------------RESUMO: PORTUGUÊS------------------------
%--------------main.tex, v1.0.0, JonathanTSilva----------------
%==============================================================

\setlength{\absparsep}{18pt} % ajusta o espaçamento dos parágrafos do resumo
\begin{resumo}
    Este é o resumo em português
    
    \textbf{Palavras-chaves}: latex; abntex; editoração de texto;
\end{resumo}

%==============================================================
%----------------------RESUMO: INGLÊS--------------------------
%--------------main.tex, v1.0.0, JonathanTSilva----------------
%==============================================================

\begin{resumo}[Abstract]

    \begin{otherlanguage*}{english}
        This is the english abstract.
        
        \vspace{\onelineskip}
        
        \noindent 
        \textbf{Keywords}: latex; abntex; text editoration;
    \end{otherlanguage*}
    
\end{resumo}
    
    %-----------------------LISTAS----------------------------|
    %-ILUSTRAÇÕES---------------------------------------------|
    \pdfbookmark[0]{\listfigurename}{lof}
    \listoffigures*
    \cleardoublepage
    
    %-TABELAS-------------------------------------------------|
    \pdfbookmark[0]{\listtablename}{lot}
    \listoftables*
    \cleardoublepage
    
    %-SIGLAS E ABREVIATURAS-----------------------------------|
    %==============================================================
%---------------------------SIGLAS-----------------------------
%--------------main.tex, v1.0.0, JonathanTSilva----------------
%==============================================================

\begin{siglas}

    \item[ABNT] Associação Brasileira de Normas Técnicas
    
    \item[abnTeX] ABsurdas Normas para TeX
    
\end{siglas}
    
    %-SÍMBOLOS------------------------------------------------|
    %==============================================================
%--------------------------SÍMBOLOS----------------------------
%--------------main.tex, v1.0.0, JonathanTSilva----------------
%==============================================================

\begin{simbolos}

    \item[$ \Gamma $] Letra grega Gama
    
    \item[$ \Lambda $] Lambda
    
    \item[$ \zeta $] Letra grega minúscula zeta
    
    \item[$ \in $] Pertence
    
\end{simbolos}
    
    %-----------------------SUMÁRIO---------------------------|
    \pdfbookmark[0]{\contentsname}{toc}
    \tableofcontents*
    \cleardoublepage
    
    %__________________________________________________________
    %------------------ELEMENTOS TEXTUAIS----------------------
    
    \textual
    
    %------------------PARTE: PREPARAÇÃO----------------------|
    %\part{Preparação da pesquisa}
    
    %-----------------CAPÍTULO: INTRODUÇÃO--------------------|
    %==============================================================
%------------------------INTRODUÇÃO----------------------------
%--------------main.tex, v1.0.0, JonathanTSilva----------------
%==============================================================

\chapter{Introdução}\label{cap:introducao}

%==============================================================
%-------------Templates de inserção de imagens-----------------
%==============================================================

%Imagem com Fonte com nota no rodapé (Disponível em: ... Acesso em:)
\begin{comment}
    \begin{figure}[H]
        \caption{\label{fig:AQUI VAI O ATALHO DA FIGURA}[AQUI VAI O CAPTION DA FIGURA]}
        \begin{center}
            \includegraphics[width=[AQUI VAI A PROPORÇAO DA FIGURA]\textwidth]{[AQUI VAI O DIRETÓRIO DA FIGURA]}
        \end{center}
        \legend{Fonte: [AQUI VAI A FONTE DA FIGURA].\protect\footnotemark}
    \end{figure}
    
    \footnotetext{Disponível em: \url{AQUI VAI O URL DA FIGURA}. Acesso em: [AQUI VAI A DATA DE ACESSO - SEMPRE COM O MES ABREVIADO: 26 jun. 2021].}
\end{comment}

%Imagem sem fonte com rodapé
\begin{comment}
    \begin{figure}[H]
        \caption{\label{fig:AQUI VAI O ATALHO DA FIGURA}[AQUI VAI O CAPTION DA FIGURA]}
        \begin{center}
            \includegraphics[width=[AQUI VAI A PROPORÇAO DA FIGURA]\textwidth]{[AQUI VAI O DIRETÓRIO DA FIGURA]}
        \end{center}
        \legend{Fonte: [AQUI VAI A FONTE DA FIGURA].}
    \end{figure}
\end{comment}

%Imagem com fonte igual de livro
\begin{comment}
    \begin{figure}[H]
        \caption{\label{fig:[AQUI VAI O ATALHO DA FIGURA]}[AQUI VAI O CAPTION DA FIGURA]}
        \begin{center}
            \includegraphics[width=[AQUI VAI A PROPORÇAO DA FIGURA]\textwidth]{[AQUI VAI O DIRETÓRIO DA FIGURA]}
        \end{center}
        \legend{Fonte: \cite{CITAÇÃO}}
    \end{figure}
\end{comment}

%==============================================================
%-------------Templates de inserção de tabelas-----------------
%==============================================================

\begin{comment}
    \begin{table}[H]
        \ABNTEXfontereduzida
        \centering
        \caption{\label{tab:[AQUI VAI O ATALHO DA TABELA]} [AQUI VAI O CAPTION DA TABELA]}
        \begin{tabular}{c c c}
            \toprule
            \thead{Item} & \thead{Nome}       & \thead{Descrição} \\
            \toprule
            1     & Controlador               & PLC S71200 CPU1214C \\
            2     & Relé Inteligente          & Simocode pro V PN \\
            3     & Transformador de Corrente & Simocode 3UF7 \\
            4     & Switch                    & SCALANCE TAP104 \\
            5     & Bomba                     & Hydrobloc P500\\
            \bottomrule
        \end{tabular}
        \legend{Fonte: [AQUI VAI A FONTE DA FIGURA]}
    \end{table}
\end{comment}
    
    %-------------CAPÍTULO: REFERENCIAL TEÓRICO---------------|
    %==============================================================
%--------------------REFERENCIAL TEÓRICO-----------------------
%--------------main.tex, v1.0.0, JonathanTSilva----------------
%==============================================================

\chapter{Referencial Teórico}\label{cap:referencial}

    
    %----------------CAPÍTULO: METODOLOGIA--------------------|
    %==============================================================
%------------------------METODOLOGIA---------------------------
%--------------main.tex, v1.0.0, JonathanTSilva----------------
%==============================================================

\chapter{Metodologia}\label{cap:metodologia}

    
    %-----------------CAPÍTULO: RESULTADOS--------------------|
    %==============================================================
%-------------------RESULTADOS E DISCUSSÕES--------------------
%--------------main.tex, v1.0.0, JonathanTSilva----------------
%==============================================================

\chapter{Resultados e Discussões}\label{cap:resultados}

    
    %-----------------CAPÍTULO: CONCLUSÃO--------------------|
    %==============================================================
%-------------------------CONCLUSÃO----------------------------
%--------------main.tex, v1.0.0, JonathanTSilva----------------
%==============================================================

\chapter{Conclusão}\label{cap:conclusao}

    
    %-----------------CAPÍTULO: EXEMPLO1----------------------|
    %\include{Textual/capExemplo1} % Capítulo com exemplos de comandos inseridos de arquivo externo
    
    %-----------------CAPÍTULO: EXEMPLO2----------------------|
    %\include{Textual/capExemplo2} % Capítulo com exemplos de comandos inseridos de arquivo externo
    
    \phantompart % Finaliza a parte no bookmark do PDF para que se inicie o bookmark na raiz e adiciona espaço de parte no Sumário
    
    %__________________________________________________________
    %----------------ELEMENTOS PÓS-TEXTUAIS--------------------
    
    \postextual
    
    %---------------------BIBLIOGRAFIA-------------------------|
    \bibliography{Postextual/references}
    
    %----------------------GLOSSÁRIO---------------------------|
    %\glossary % Consulte o manual da classe abnTeX2 para orientações sobre o glossário.

    %----------------------APÊNDICES---------------------------|
    %==============================================================
%-------------------------APÊNDICES----------------------------
%--------------main.tex, v1.0.0, JonathanTSilva----------------
%==============================================================

\begin{apendicesenv}
    
    \partapendices % Imprime uma página indicando o início dos apêndices
    
    \chapter{Exemplo1}
    
\end{apendicesenv}


    
    %------------------------ANEXOS----------------------------|
    %==============================================================
%--------------------------ANEXOS------------------------------
%--------------main.tex, v1.0.0, JonathanTSilva----------------
%==============================================================

\begin{anexosenv}
    
    \partanexos % Imprime uma página indicando o início dos anexos
    
    \chapter{Exemplo1}
    

\end{anexosenv}
    
    %-------------------ÍNDICE REMISSIVO-----------------------|
    \phantompart
    \printindex
    
\end{document}


